\section[PHT3D Exercise 6, Zinc transport at the Cape Cod site]{PHT3D Exercise 6, Zinc transport at the Cape Cod site}

In this exercise we will construct a model that provides 
a good example for the consideration of surface complexation reactions in PHT3D.
The example is based on a modelling study for the Cape Code site, 
The full details of this study were reported in \cite{Kent20003411}.    
The Cape Cod research site is near a wastewater-treatment facility at 
the Massachusetts Military Reservation (MMR). A treated-sewage plume originated 
from the facility's infiltration beds, which were used from about 1936 to 1995. 
At the time the modelling study was undertaken the plume at the site extended 
more than 6 kilometers from the disposal site in the 
sand and gravel aquifer. The plume contains a complex mixture of phosphate, nitrate, 
metal ions (such as zinc), detergents, organic chemicals, and microbes. 
Over many years the plume and adjacent parts of the aquifer  
served as a field laboratory for a multidisciplinary team of scientists 
who investigated many fundamantal aspects 
of contaminant transport behaviour.

\subsection{Background: Groundwater Flow}

At the Cape Cod site the regional groundwater flow is generally towards 
the south, but locally the flow direction is directed from the 
sewage treatment facility towards the eastern shore of Ashumet Pond.  
The groundwater flow is nearly horizontal.
Approximately 25 cm/yr of areal recharge results
in a small vertical component to flow as well as
in the accumulation of a lens of pristine groundwater on top of 
the sewage plume, the thickness of which increases with 
distance away from the sewage 
treatment plant \citep[]{Leblanc1984,Leblanc1991895}.  
The magnitude of the hydraulic gradient 
is approximately 1.7 m/km. 
Seasonal fluctuations in the flow system result in small 
variations (approximately 0.2 m/km) in the magnitude
of the gradient as well as in small annual fluctuation 
in flow direction.

\subsection{Background: Hydrogeochemistry}

Contrasting chemical conditions between sewage-contaminated
and ambient groundwater gave rise to longitudinal and 
vertical gradients in groundwater chemistry. 
The ambient groundwater surrounding the sewage plume had 
high concentrations of dissolved oxygen, pH values 
in the range 5.3-5.8, and low concentrations of dissolved salts 
\citep[]{Leblanc1991895,Hess1996,Kent1999293}.
%[LeBlanc  et  al., 1991; Hess  et al., 1996; Kent and Maeder, 1999].
Sewage-contaminated groundwater under the disposal beds had high 
concentrations of nitrate.
Concentrations of dissolved oxygen successively decreased with depth 
\citep[]{Smith1991285,Smith19911815}.
%[Smith et al., 1991a].  
Biodegradation of organic matter in conjunction with solute transport 
away from the disposal beds resulted in the depletion of dissolved oxygen 
throughout the plume and the depletion of nitrate, followed by 
the accumulation of dissolved Fe(II), in the core of the plume 
%[Smith et al., 1991b; Kent  et  al., 1994]
\citep[]{Smith1991285,Kent19941099}. 
 
These biodegradation reactions influenced the pH  
within the aquifer zone affected by the sewage plume 
%[Abrams et al., 1998; Abrams and Loague, 2000b].
\citep[]{Abrams19981531,Abrams20002015}.
The principal chemical process in terms of controlling Zn transport is 
assumed to be Zn adsorption onto the sediments.
A single reaction with the stoichiometry could already adequately 
describe the pH effect observed in laboratory batch studies of
Zn adsorption onto the sediments \citep[]{Davis19982820}:

\begin{equation}
Zn^{2+}  +  >SOH    =  >SOZn^+  + H^+,  log  ^SQ_{SOZn}  
\label{eq:zn_single_site}
\end{equation} 

Adding a second surface site with the same reaction stoichiometry results in an  
SCM that can fit both the pH and Zn-concentration 
effect \citep[]{Davis19982820}.
This so-called two-site approach was one of the various
numerical implementations tested by \cite{Kent20003411}.
This two-site model will be used in the present exercise.
In the two-site SCM a small percentage of the total 
adsorption sites are considered to have a much greater affinity 
for adsorption ("strong sites," $>SsOH$) than the rest of the 
sites ("weak  sites," $>SwOH$): 

\begin{equation}
Zn^{2+}  +  >S_sOH    =  >S_sOZn^+  + H^+,  log  ^SQ_{S_sOZn}  
\label{eq:zn_strong_sites}
\end{equation} 

\begin{equation}
Zn^{2+}  +  >S_wOH  =  >S_wOZn^+  + H^+,  log  ^SQ_{S_wOZn}   \\
\label{eq:zn_weak_sites}
\end{equation} 

Before proceeding, it is worthwhile to consider the standard state for sorbed species 
in surface complexation models. 
The same standard state used for aqueous species - either 1 mole of surface complex \textbackslash L of solution (M)or \textbackslash kilogram water (m) - is often used for surface complexes as well. For example, the comprehensive compilation of equilibrium constants for sorption on hydrous ferric oxide by \cite{Dzombak1990} adopts a standard state of 1 M for sorbed species.  
One should be aware, however, there is a pitfall associated with this choice of 
standard state %\citet{appelo_postma_1999})
\citep[e.g.,][]{appelo_postma_1999}. 
The pitfall can be illustrated by considering a batch reactor with a suspension in sorptive equilibrium. 
In this thought experiment, we allow all of the solids to settle to the bottom of 
the reactor once sorptive equilibrium has been achieved, then we remove half of the solution without removing any solid.
Sorbed concentrations expressed in moles/kg water have now increased by a factor of two because the volume of solution (hence, mass of water) 
has decreased by a factor of two with no change in the amount of the solid.  
For surface complexation reactions such as those described the Eqn. \ref{eq:zn_strong_sites} 
and Eqn. \ref{eq:zn_weak_sites}, 
the concentration of sorbed Zn doubles and the concentration of free surface functional groups double,
so the factor of two cancels out in the mass action expression.  
For this type of surface complex - 
called ''monodentate'' (one surface site per sorbed species) -
the choice of standard state for sorbed species does not affect the value of the equilibrium constant. 
However, for ''bidentate'' surface complexes, where there are two surface functional groups 
bound to a single sorbing species, the situation is more complicated.  
Consider formation of a bidentate surface complexation analogous 
to Eqn. \ref{eq:zn_strong_sites}:

\begin{equation}
Zn^{2+}  +  2>S_sOH    =  (>SsO)2Zn + 2H^+  \\
\label{eq:zn_2}
\end{equation}

The mass action expression can be written with the concentration 
of the Zn surface complex raised to
the first power in the numerator and concentration of free surface 
sites squared in the denominator.  
In this mass action expression, the factor of two does not cancel out.  
This observation can be generalized by stating that equilibrium constants describing sorption where multidentate surface complexes are involved are a function of the solid to liquid ratio if a standard
state of 1 M or 1 m is adopted for sorbed species.  
There has been much discussion of this in the literature %(e.g.,\citet{wang_giammar_2013})
\citep[e.g.,][]{wang_giammar_2013}.  
To avoid this pitfall, PHREEQC adopts a standard state for sorbed species whereby the fraction 
of sites occupied by a given sorbed species equals 1.
In reaction \ref{eq:zn_2}, the relationship between sorbed Zn 
in units of moles \textbackslash kg water and 
sorbed Zn in units of mole fraction occupied 
sites is 
\begin{equation}
XZn = 2[(>SsO)2Zn]/[SsOHT],  \\
\label{eq:xzn}
\end{equation}
where square brackets refer to concentrations in moles/kg water and SsOHT is the total 
concentration of strong sorption sites.  
The factor of two results from the fact that the surface complex covers two sites.  
Therefore, when using equilibrium constants for multidentate sorbed species taken 
from the literature in a transport simulation, it is important to insure 
that either the equilibrium constant uses mole fraction of sites occupied 
by any given sorbed species = 1 or is properly converted to that standard state.  
Formulas for making such conversions are presented 
in various papers 
%\citep[see][]{jon90} 	    -->    	(see Jones et al., 1990)
\citep[e.g.,][]{wang_giammar_2013}.

\par With the sorption and thus the mobility of Zn being strongly 
depedent on pH, it is important to accurately capture the pH-controlling 
processes. For the vertical transect model that will be constructed
in this exercise the evolution of the pH will mostly be 
affected by (i) the transport and mixing of two different water 
types, i.e., of the (upstream) water composition
and the slightly more acidic (pH = 5.6) groundwater recharge water and (ii)
the pH increase that is induced by the degradation of organic matter
in conjunction with sulfate and iron reduction. The former causes
a pH gradient that is the result of physical transport processes. 
The latter causes inside the plume-affected aquifer zone
a pH increase towards 6.5, and thus above the 
pH of 6.0 that characterises the uncontaminated background water. 
For this exercise the organic matter degradation 
is represented in a simplistic way by a single kinetically 
controlled degradation reaction of DOC in conjunction
with the reduction of two electron acceptors sulfate and goethite.
It is assumed that nitrate is already depleted at the model's 
upstream boundary and thus not considered.
The rate $r_{DOC}$ of the DOC mineralisation is computed from:

\begin{equation}
r_{DOC} = \frac{C_{DOC}}{K_{DOC} + C_{DOC}} \sum_{i=1,nr_{EA}}{k_{EA_i} \frac{C_{EA_i}}{K_{EA_i} + C_{EA_i}}} 
\label{eq:doc_deg}
\end{equation}

where $C_{DOC}$ is the DOC concentration, $C_{EA_i}$ is the concentration of the $i^{th}$
electron acceptor (in the present case either sulfate or goethite) while $K_{DOC}$ and $K_{EA_i}$
are half saturation constants for DOC and the $i^{th}$
electron acceptor, respectively.

\par In reality the composition of the sewage plume and the 
reactions are obviously far more complex, as discussed in 
detail by \cite{Kent20003411}.
Nevertheless, the inclusion of this reaction allows to replicate
some of the key features observed at the site, and, most importantly,
allows to mimic the observed pH patterns relatively well.
Therefore the model constructed in this exercise also reproduces
the major Zn transport patterns relatively well, i.e., the rates 
of Zn movement and the vertical differences in the mobility, as induced
by the spatial variability of the pH.  

\begin{table}[hb]
\caption{Flow and solute transport parameters used for the setup of PHT3D Exercise 5}
\begin{center}
\begin{tabular}{lc} 
\toprule
Parameter                                            &         Value \\
\midrule
Flow simulation type                                 &  steady state  \\
Simulation time (days)                               &  21000         \\
Time step (days)                                     &  60            \\
Model \  extent\ column\ direction\ $ (m)$           &  400           \\
Model \ top\ elevation\ $(m)$                        &   16           \\
Model \ bottom\ elevation\ $(m)$                     &   0            \\
Horizontal hydraulic conductivity\ $(m/day)$         &   110          \\ 
Porosity, $\theta$                                   &   0.39         \\
Piezometric head downstream boundary\ $(m) $         &  14            \\ 
Longitudinal dispersivity, $\alpha_L$   $(m)$        &   1.0          \\
Horizontal transverse dispersivity, $\alpha_{TH}$    &   0.1          \\
Vertical transverse dispersivity,   $\alpha_{TV}$    &   0.01         \\
Effective diffusion coefficient, D$^{\star}$         &   0            \\
\bottomrule

\end{tabular}
\label{cape_cod_table_1}
\end{center}
\end{table}

\subsection{Setting up the flow model}

To set up the flow model proceed as follows:

\begin{itemize}

\item In ORTI, create a new model and save it in a new and separate directory (folder).

\item Select {\color{blue}Parameters} $\rightarrow$ {\color{blue}1.Model} 
$\rightarrow$ {\color{blue}Model} to specify the model dimension {\color{blue}('Xsection')}
and type {\color{blue}('unconfined')}

\item Select {\color{blue}Parameters} $\rightarrow$ 
{\color{blue}1.Model} $\rightarrow$ 
{\color{blue}Domain} to specify the domain as described in Table 
\ref{cape_cod_table_1}. 
For the first simulation select a spacing of 20 m in 
x direction and 1 m in y direction

\item Select {\color{blue}Parameters} $\rightarrow$ {\color{blue}Time} 
$\rightarrow$ and set the {\color{blue}Total Simulation Time} 
to 21000 days. Then define the step size to 60 days

\item Set the boundary conditions for heads by selecting 
{\color{blue}Spatial Attributes} $\rightarrow$ {\color{blue}MODFLOW} 
$\rightarrow$ {\color{blue}bas.3 Boundary Conditions} corresponding to
IBOUND (Modflow). 
Use the {\color{blue}Zone Tool} to create a zone at the 
outflow end that extends vertically between z = 0 m and z = 13.9 m.

\item To fix the hydraulic heads at the boundaries, use 
{\color{blue}Spatial Attributes} $\rightarrow$ 
{\color{blue}MODFLOW} $\rightarrow$ {\color{blue}bas.5 Initial Heads} and provide 
a value of 16 in the dialog box. This value will be used everywhere as a 
starting (initial) head and will be kept at this value at all locations
where a fixed head boundary condition was defined in the previous step. 
In order to fix the heads at the outflow boundary, 
draw a line at the same location as the outflow BC 
and attribute a head value of 14 m.

\item The inflow boundary is implemented by a single well that vertically extends 
between 0 and 15 m. The total injection rate of this well should be set to 2.16 $m^3/d$.

\item Select {\color{blue}Parameters} $\rightarrow$ {\color{blue}2.Flow} 
$\rightarrow$ {\color{blue}Parameters} 
and use this dialog to set several parametters. 
In {\color{blue}lpf.2} select {\color{blue}Convertible} 
for the {\color{blue}type of layer}. 
Then use {\color{blue}lpf.8} to specify the hydraulic conductivity.

\item To fix recharge, use {\color{blue}Spatial Attributes} $\rightarrow$ 
{\color{blue}MODFLOW} $\rightarrow$ {\color{blue}rch.2 Recharge}. 
Create a zone along the top cells and define a value of 0.0014 in the dialog box.

\item To define the porosity select {\color{blue}Spatial Attributes} 
$\rightarrow$ {\color{blue}MT3DMS} $\rightarrow$ {\color{blue}btn.11} and 
set {\color{blue}Value} to 0.39. 

\end{itemize}

Regularely save the changes to avoid loss of your data input.

\subsection{Running MODFLOW}

You have now completed the set up the flow model and are ready to run MODFLOW.
This step will provide PHT3D with the flow field that will be used to simulate 
reactive transport.

\begin{itemize}

\item Select {\color{blue}Parameters} $\rightarrow$ {\color{blue}2.Flow} $\rightarrow$ {\color{blue}Write}

\item Select {\color{blue}Parameters} $\rightarrow$ {\color{blue}2.Flow} $\rightarrow$ {\color{blue}Run}

\item When the calculations have finished, a new window will appear. I will show the last line(s) of the MODFLOW output file. Make sure MODFLOW terminated without errors.

\end{itemize}

\subsection{Visualising hydraulic heads}

\begin{itemize}

\item Visualise results by selecting {\color{blue}Results} $\rightarrow$ {\color{blue}2.Flow} $\rightarrow$ {\color{blue}Head}. The time step can be changed
in {\color{blue}Results} $\rightarrow$ {\color{blue}1.Model} $\rightarrow$ {\color{blue}Time}.

\item Assess if the head distribution is plausible.

\end{itemize}

\subsection{Chemistry}

In this section we will define (i) the initial hydrochemical 
composition (ii) the solutions that represent the sewage
plume composition at the upstream model boundary
and (iii) the groundwater recharge composition. 
Proceed with the following steps:

\begin{itemize}

\item Import the database from the folder that contains the Supporting Information for this exercise


\item The species needed to be activated are $Cl$, $C(4)$, 
$Fe(2)$ and $Fe(3)$, $Na$, $O(0)$, $S(-2)$ and $S(6)$, $Orgc$, $Zn$ and, as always, $pH$ and $pe$. 
As described in the previous PHT3D Examples, the activation of these species/componenets 
and the definition of the solution compositions are done in the {\color{blue}Solution Tab}. 
All required concentrations are listed in Table \ref{scm_solution_table}.

\item In is assumed that goethite is present at a (background) concentration 0.1 $mol/L_b$. 
FeS(ppt) must also be ticked. It is initially not present and thus the concentration
is left at the default value of 0.

\item The allocation of the solutions to define initial and boundary conditions
can be done via {\color{blue}ph.4 Source / Sink} zones (line). 
To allocate the correct solutions at the inflow boundary allocate 
from 14 m to 16 m $Solution\ 2$, from 3 m to 14 m 
allocate $Solution\ 1$ and between 0 m and 3 m allocate the background solution,
i.e., $Solution\ 0$.

\item The groundwater recharge composition is entered via {\color{blue}ph.5} as a line along 
the top grid cells. The recharge water composition is defined by $Solution\ 3$.
However, make sure that the recharge does not extend to the cell at the inflow end of the grid.   

\item Under the {\color{blue}Rates Tab} specify the degradation rate 
parameters for Orgc (0,1e-11, 1e-12, 1e-12,1e-17,1e-17 respectively) and also 
enter the stoichiometry of the mineralisation 
reaction. The formula is 
Orgc -1.0 CH2O 1.0.
%Orgc -1.0 CH2O(NH3)0.15 1.0.

\item Activate both available types of surface sites, i.e., $S\_s$ and $S\_w$.
Define 1.59e-5 $mol/L_b$ and .0018 $mol/L_b$ as their respective background concentrations.


\end{itemize}

\subsection[Running PHT3D]{Running PHT3D}

With the model setup completed run PHT3D and visualise the results via contour plots.


\begin{table}[hb]
\caption{Chemical composition for PHT3D Exercise 5}
\begin{center}
\begin{tabular}{lcccc} 
\toprule
Parameter  & Background & Solution 1 & Solution 2 & Solution 3\\
\midrule
C(4)      &   1e-4 &   1e-4  &   1e-4 &   1e-4\\
Cl      &   0 &   0  &  0  &   1e-4 \\
Fe(2)      &   0 &   0  &  0  &   0 \\
Fe(3)      &   0 &   0  &  0  &   0 \\
Na      &   2.3e-4 &   2.3e-4 &   2.3e-4 &   3.13e-4 \\
O(0)      &   5e-4 &   0 &   5e-4 &   5e-4\\
Orgc      &   0 &   1e-4 &   0  &  0 \\
S(6)      &   1e-4 &   1e-4  &   1e-4 &   1e-4\\
S(-2)      &   0 &   0  &  0  &   0 \\
Zn      &   0 &   8e-6  &  8e-6  &   0 \\
pH      &   6 &   6 &  6  &   5.6 \\
pe      &   14.6 &   4 &  14.6  &   14.2 \\
\bottomrule

\end{tabular}
\label{scm_solution_table}
\end{center}
\end{table}

